\documentclass[11pt]{article}
\usepackage[utf8]{inputenc}
\usepackage[T1]{fontenc}
\usepackage[margin=1in]{geometry}
\usepackage{amsmath,amssymb}
\usepackage[hidelinks]{hyperref}
\setlength{\parindent}{0pt}
\setlength{\parskip}{0.6em}

\title{Response to the Mathematical Review\\
\large Informativity and Identifiability for Identification\\
of Networks of Dynamical Systems}
\author{Anders Hansson, Jo\~ao Victor Galv\~ao da Mata,\\Martin S.~Andersen}
\date{September 14, 2026}

\begin{document}
\maketitle

This response records the changes present in the revised manuscript as of the
date above. It responds to the supplied \emph{Detailed Mathematical and Logical
Review}, dated September 3, 2026. References below are to the sections of that
internal technical review; comments are paraphrased rather than quoted.
We distinguish corrections relevant to the stated results from suggestions for
additional results or alternative presentations. Not adopting the latter does
not constitute an unresolved objection to the results considered here.

\section*{General response}

We thank the reviewer for the detailed mathematical assessment and the explicit
examples illustrating the distinctions between actual and generic rank,
admissible sets and their closures, and generic and exceptional fibers.
We have revised the manuscript to clarify these distinctions while retaining its
central framework and its focus on Gr\"obner basis calculations.

The main changes are a generic qualification of the informativity theorem,
an explicit row--column matching convention for the graph construction,
a more complete specification of the admissible denominator domain, and an
algebraic formulation of generic local network identifiability over
$K=\mathbb C(z)$. We have also corrected or clarified the five examples and
aligned the abstract, introduction, and conclusions with the revised scope.
The responses below explain how the relevant mathematical objections have been
addressed by corrections, explicit qualifications, or clarification of existing
assumptions. Where a suggestion is not relevant to the scope of the claimed
result, we say so and explain why. Separate editorial and reproducibility
matters are distinguished from these mathematical responses at the end.

\section{Predictor algebra and spectral preliminaries}
\textbf{Review Sections 2.1--2.2.}
The reviewer accepts the predictor algebra subject to the necessary assumptions,
but requests explicit rational invertibility and innovation assumptions and a
full spectral-integral argument.

\textbf{Response.}
The paper assumes the existence of the innovation representation and a stable
predictor; it does not attempt to establish existence of that representation
from the stated distributional properties of the innovations alone. Such an
existence theorem is not relevant to the conditional results studied here.
The required rational inverses in the displayed predictor formulas are
understood as part of that representation, not as consequences of stability
alone. In particular, wherever the formulas are defined,
$I-W_y=(I+G_o)^{-1}$. We retain the existing predictor algebra rather than add a
separate derivation of the innovation model.

The request concerning Parseval's formula is relevant as a clarification of
the reasoning, but does not require a different informativity argument.
For completeness, under the second-order stationarity assumptions implicit in
the spectral calculation, write $d_k=\Delta W(q)z_k$. Then
\[
\mathbb E[d_kd_k^*]=\frac{1}{2\pi}\int_{-\pi}^{\pi}
\Delta W(e^{i\omega})\Phi_z(\omega)\Delta W(e^{i\omega})^*\,d\omega.
\]
The integrand is positive semidefinite, so a zero integral is equivalent to a
zero integrand almost everywhere. Positive definiteness of $\Phi_z$ then
implies $\Delta W=0$ almost everywhere. This is the step abbreviated in the
manuscript. For real-coefficient transfer matrices, the manuscript's notation
$\Delta W(e^{-i\omega})^T$ equals $\Delta W(e^{i\omega})^*$; it is not a
transpose evaluated at the same frequency. The revised proof of Theorem~1
also explicitly uses the true-network spectral identity with a conjugate
transpose. The expanded integral is supplied here as an explanation, not
reported as an addition to the manuscript.

\section{Actual versus generic informativity}
\textbf{Review Sections 3.1--3.3.}
Full generic row rank does not imply full row rank for every true network, since
the true network can belong to an exceptional rank-deficient set.

\textbf{Response.}
We agree and have revised Theorem~1 and its proof. The theorem now states that,
under the standing assumptions, positive definiteness of the external-signal
spectrum almost everywhere and full generic row rank of $\Pi$ imply
informativity for almost all admissible choices of the true network's free
transfer functions. The proof uses
\[
\Phi_z(\omega)=\Pi(e^{i\omega};\theta_0)\Phi_u(\omega)
\Pi(e^{i\omega};\theta_0)^*.
\]
The text following the theorem explicitly allows exceptional true networks,
including networks for which cancellations reduce the rank. It also clarifies
that the generic qualification concerns the true network generating the data,
not the pairs of candidate predictors in the definition of informativity.
Generic rank is understood relative to the model class, including its known
transfer functions and constraints.

The fixed-system sufficient condition remains in the discussion preceding the
theorem. Thus, we retain that condition and give a qualified generic theorem,
rather than introducing a separate fixed-system theorem and corollary.
The corresponding generic qualifications have been propagated to the numerical
examples and the summary sections.

\section{Reduced predictors and nuisance columns}
\textbf{Review Section 3.4 and the predictor aspect of Section 6.6.}
An unknown predictor column cannot be discarded merely because it is not of
interest: its variation can cancel a change in another column.

\textbf{Response.}
We have rewritten Remark~1. Measurement rows may be omitted when their
corresponding predictor columns are fixed across the candidate models, but not
merely because unknown columns are of no interest. The remark also distinguishes
this operation from removing external-input columns when the associated inputs
are set to zero.

We additionally explain that known relations between predictor columns can
permit a reduced vector of known combinations of measured signals.
Example~2 now uses this construction explicitly instead of relying on deletion
of an unknown predictor column. This addresses the cancellation issue for the
reduction actually used in the paper. A general projected-fiber or
projected-kernel theorem for selected modules with arbitrary nuisance parameters
is not relevant to establishing this particular reduction and would extend the
scope of the paper. We therefore do not add such a theorem or claim that
arbitrary nuisance directions can be ignored.

\section{Graph-theoretic generic rank}
\textbf{Review Sections 4.1--4.4.}
The directed graph requires a justified pairing of the equation rows and node
columns of $P$. Moreover, structural matching arguments cannot in general treat
fixed or algebraically dependent entries as independent free entries.

\textbf{Response.}
We retain the block-rank identity and have added an explicit convention before
the graph construction: the diagonal entries of $P$ are structurally nonzero.
For a generically invertible $P$, a row permutation can select such a diagonal;
the same permutation is applied to $Q$. This leaves $RP^{-1}Q+S$ unchanged.
The matrices used in the subsequent graph construction are explicitly the
matrices after this permutation. Thus, the same-index row--column pairing is
made after selecting a structural matching, not imposed on an arbitrary
unpermuted $P$.

We have retained the directed-path presentation rather than replacing it with a
separate general bipartite-matching theorem. The accompanying remark states that
the graph proof requires the nonzero entries of $M$ to be free variables.
The introduction and conclusions now qualify the graph characterization as a
free-entry result and distinguish it from the algebraic treatment of known
transfer functions and constraints. The description of the exceptional
rank-deficient set has also been qualified as being contained in a proper
affine variety.

\section{Admissible domains and the Gr\"obner basis rank test}
\textbf{Review Sections 5.1--5.4.}
The denominator exclusion must include well-posedness and all rational
parameterization or constraint denominators. Elimination describes a closure,
and minors of order $k$ detect rank below $k$, not necessarily below the generic
rank.

\textbf{Response: changes made.}
In the rank construction, we now use
\[
D=\det(P)\prod_i D_i,
\]
with additional denominator factors from rational constraints or
parameterizations included when required. The equation $1-tD=0$ enforces
$D\ne0$, and the text states that the subsequent rank statements and solution
sets are understood on this admissible domain. We have also removed the
misleading removable-pole footnote. In the identifiability calculation, the
corresponding denominator product is enlarged to include the polynomial factors
excluded from the admissible domain.

We have also corrected item~3 of Theorem~3: the polynomials independent of $t$,
together with $D\ne0$, define the locus $\operatorname{rank}\Pi<k$.
Choosing $k=\operatorname{rank}_g\Pi$ gives the exceptional set below generic
rank. A sentence following the theorem explicitly distinguishes this admissible
set from the algebraic closure defined by the elimination ideal without the
restriction $D\ne0$.

The constraint remark now requires denominators of rational equations to be
cleared and included in $D$. The zero-ideal test is interpreted in the coordinate
ring of the relevant reduced irreducible component. Reducible constraint sets
are treated componentwise, considering only components meeting $D\ne0$.
These revisions address the rank-threshold, admissibility, closure, and
constraint qualifications raised in the review. They retain the existing
elimination construction rather than introduce a different algorithm.

\section{Generic local network identifiability}
\textbf{Review Sections 6.1--6.5.}
The image of the rational map need not be an affine variety, exceptional fibers
can have larger dimension, and symbolic rational functions must be distinguished
from transfer-function values at a single frequency.

\textbf{Response.}
We have revised the definitions, Theorem~4, its proof, and the computational
remarks as follows.
\begin{enumerate}
\item We keep $z$ symbolic and define
$F(X)=R_yP^{-1}Q_r+S_{yr}$ over $K=\mathbb C(z)$.
Known rational transfer functions enter as coefficients. Algebraic varieties
and fibers are considered over an algebraic closure of $K$; rational
transfer-function networks correspond to points with coordinates in $K$.
Equality under $F$ is equality of rational functions, not equality at a fixed
frequency.
\item The admissible domain $V_o$ excludes singular $P$ and points where the
expressions defining $F$ are undefined. We explicitly state that
$V_c=F(V_o)$ need not be affine and use its dimension to mean the dimension of
its algebraic closure.
\item We assume irreducibility of $V_o$ and define generic local network
identifiability in the algebraic sense through zero-dimensional generic
fibers. Theorem~4 now concerns generic image points, not every admissible
$G_c$. Its proof applies the generic fiber-dimension formula.
\item With constraints, we use the actual dimension of the constrained domain,
not automatically the number of listed variables. Reducible cases are treated
componentwise, together with the closure of each component's image.
\item We explain that the zero-dimensional fibers give finitely many isolated
algebraic solutions. Such solutions need not all be rational transfer functions;
properness and stability must still be checked when interpreting alternatives
as admissible networks. The former Euclidean-neighborhood formulation has been
replaced by this algebraic interpretation, with the distinction from the
single-frequency formulation in the cited work made explicit.
\end{enumerate}

We also acknowledge the Jacobian characterization in characteristic zero.
In the unconstrained case, the Jacobian of $F$ must have generic rank $k$ over
$K(X)$. With constraints, the differential is restricted to the tangent space
at a generic smooth point and must have rank $\dim(V_o)$. We retain the
Gr\"obner basis formulation as the computational focus of the paper.

\section{Subnetwork constructions}
\textbf{Review Section 7.}
The boundary signals used as inputs must actually be measured or otherwise
available; an unknown filter does not permit a filtered boundary signal to be
computed from an unfiltered measurement alone.

\textbf{Response.}
We have retained the constructions and measurement assumptions of Section~VI,
with the notation and matrix-sign corrections described below. The section
explicitly distinguishes the use of
measured node signals from the alternative in which filtered boundary signals
are measured. The latter is a measurement assumption, not a claim that an
unknown filter can be applied numerically to reconstruct those signals.
Unknown transfer functions retained in the subnetwork model remain part of
its identification problem. The section is read with the revised generic-rank
and algebraic-identifiability qualifications. The objection that an unknown
boundary filter cannot be used to compute a signal is therefore not applicable
to the stated alternative: the filtered signal itself is assumed measured.
No additional change to Section~VI is required on that account.

\section{Examples}
\subsection*{Example 1: Review Section 8.1}
We have corrected the missing noise-filter factors in $\Pi$ and the sign of its
determinant, which is now
\[
\det\Pi=\frac{H^1H^2}{1-G^1G^2}.
\]
The nonvanishing conditions on the noise filters and the standing well-posedness,
stability, and spectral assumptions are stated. The discussion of the map
$(G^1,Q^1)\mapsto(G^1,G^1Q^1)$ distinguishes its generic inverse from the
exceptional fiber at $(0,0)$, where $Q^1$ is arbitrary. This also illustrates why
the revised Theorem~4 is a generic-fiber statement rather than a statement about
every admissible image point.

\subsection*{Example 2: Review Section 8.2}
With known $Q^1$, we now form the measured combination
$\xi_k=w_k^3+Q^1(q)r_k$ and write
\[
\hat w_k^1=W_y(q)w_k^1+V(q)\xi_k,
\qquad G^1=(1-W_y)^{-1}V.
\]
The original predictor columns are related as $(W_y,V,VQ^1)$, so this is a
reduced predictor based on known relations, not deletion of an unknown column.
The example establishes the required rank conditions for the input pair
$(w^1,\xi)$ using the determinants
\[
-\frac{H^1(Q^1+G^2Q^2)}{1-G^1G^2}\quad(e^2=0),
\qquad
\frac{H^1H^2}{1-G^1G^2}\quad(r=0).
\]
The conclusion is qualified by nonzero determinants, positive-definite spectra
of the active external inputs almost everywhere, and the standing assumptions.

\subsection*{Example 3: Review Section 8.3}
We now explicitly use two separate noise-free experiments, exciting $r^1$ in
one and $r^2$ in the other, with the corresponding measured outputs stated.
With the stated excitation and stability assumptions, the closed-loop transfer
functions are determined directly from these input-output relations. The text
explicitly avoids claiming unique determination of an unrestricted predictor
from a spectrally dependent noise-free input-output vector.

We retain the twofold algebraic ambiguity and clarify that the duplicated-network
construction involves equality constraints between repeated transfer functions.
The conclusion is framed as a limitation of extending the cited conjecture to
that enlarged model class, not as an unqualified counterexample for the original
class. The node-signal typographical error has also been corrected.

\subsection*{Example 4: Review Section 8.4}
We retain the reported dimension calculation and its generic local network
identifiability conclusion. The informativity statement now explicitly refers
to full generic row rank, generic informativity, the standing assumptions, and
positive definiteness of the external-signal spectrum almost everywhere.
We appreciate the independent Jacobian certificate supplied in the review,
which supports the reported conclusion. The Jacobian criterion is acknowledged
in the general theory. Adding this particular alternative certificate is not
necessary to correct the example; it is an alternative computational
presentation. We retain the Gr\"obner basis calculation because illustrating
that formulation is the purpose of the numerical example.

\subsection*{Example 5: Review Section 8.5}
We have corrected the unknown tuple to $(G^{12},Q^1,Q^2,R^1)$ and the measured
output to $y_k=(R^1(q)w_k^1,w_k^2)$, consistently with the displayed $R_y$ and
$F(X)$. The filtered boundary signal $\tilde w_k^B=G^{23}w_k^3$ remains explicitly
assumed measured.

The reduced-excitation discussion now removes the inactive input columns and
the corresponding measured-input rows from $M$, rather than leaving zero rows
in the rank test. Removing $r^3$ additionally removes the last column of $F(X)$.
The informativity conclusions are qualified generically and by the spectral
condition on the remaining active external signals. We have retained the focus
on generic local network identifiability. The proposed explicit inverse and
stronger uniqueness discussion are not relevant to establishing the local
identifiability conclusion reported in this example. We therefore do not extend
the example to that stronger question; this is a scope choice, not an unresolved
objection to the stated conclusion.

\section{Summary sections and scope of the revision}
\textbf{Review Sections 9--11.}
The abstract, introduction, and conclusions now refer to generic informativity
and full generic row rank. The algebraic sense of generic local network
identifiability is made explicit. The introductory and concluding dimension
statements refer to the admissible domain and the closure of its image. The
graph characterization is qualified by the free-entry setting, and the claim
about known transfer functions reducing computation time has been softened
from a guaranteed reduction to a possible reduction.

The irreducibility and image-relative genericity points raised in Review
Section~9 are addressed in the identifiability formulation. The relevant
mathematical corrections in the proposed replacement package have been
incorporated into our existing presentation, including the latest changes to
Theorem~3. Reorganizing the paper into the reviewer's preferred theorem package
is not necessary for those corrections and has not been adopted.

\section*{Separate editorial and reproducibility matters}
The following matters should not be confused with requests to broaden the
mathematical scope. They are recorded separately so that the response does not
claim manuscript changes or supporting material that have not been supplied.

\textbf{Notation and typographical corrections: completed.}
We have corrected the notation and typographical points raised in the review:
\begin{itemize}
\item The dimension of $r_k$ is consistently denoted by $m_r$, with
$p=p_y+m_r$; the redundant notation $p_r$ has been removed.
\item The closed-loop transfer formula now uses the output feedthrough blocks
$S_{yr}(q)$ and $S_{ye}(q)$, consistently with the partition of the measurement
equation.
\item All displayed $M$-matrices follow the convention
$M=\left[\begin{smallmatrix}P&Q\\-R&S\end{smallmatrix}\right]$.
We have corrected the node-signal measurement blocks in the subnetwork matrices
and Example~5, including the rows for measured boundary signals. The direct
feedthrough blocks retain their original signs. The underlying network and
measurement equations are unchanged.
\item The predictor difference is now written as
$\hat y_k(\theta_1)-\hat y_k(\theta_2)=\Delta W(q)z_k$, consistently with
$\Delta W=W(q;\theta_1)-W(q;\theta_2)$. This sign correction does not change
the zero-covariance condition.
\item The bibliography year 2913 has been corrected to 2013. We have also
corrected the spelling of ``partitioned'', the subnetwork replacement label
$Q_r$, the missing comma in the filtered-boundary input tuple, and the subscript
$R_{yA}$.
\end{itemize}

\textbf{Computational documentation.}
The request for code, polynomial systems, ordering conventions, software details,
and certificates is relevant to reproducibility of the reported computations.
No new archive is supplied by the present revision. This request is distinct
from requiring a particular alternative Jacobian calculation or a stronger
identifiability conclusion, neither of which is needed for the scope of the
examples. We do not claim that an unprovided archive has been added.

\end{document}